La utilidad del sistema híbrido DEB+GP no reside únicamente en su capacidad para generar predicciones, sino en la demostración cuantitativa de que dichas predicciones superan, en precisión y robustez, a los métodos de estimación estática convencionales. Esta sección establece el protocolo de validación predictiva, las métricas de regresión empleadas y el análisis de residuos orientado a la detección de eventos de anaerobiosis. Los criterios de éxito se alinean con las hipótesis H2, H3 y H4 formuladas en el \cref{chap:planteamiento_problema}: precisión predictiva del modelo mecanicista refinada por el estimador GP, y sensibilidad operativa a condiciones transitoriamente anaeróbicas.

\subsection{Coeficiente de determinación y error porcentual absoluto medio}
\label{subsec:cap4_metricas_regresion}

La validación del desempeño predictivo se estructura como una comparación contra un modelo de referencia de línea base. Este modelo estático, denominado \emph{lineal de tendencia}, ajusta una regresión lineal simple entre el tiempo transcurrido $t$ y la pendiente de emisión observada $\Delta C/\Delta t$, sin incorporar información sobre biomasa, temperatura ni humedad. Representa el nivel de predicción que un operador podría obtener mediante extrapolación empírica directa, y constituye el umbral mínimo que el sistema híbrido debe superar para justificar su complejidad adicional \parencite{weichertReviewMachineLearning2019}.

Sea $\{y_k\}_{k=1}^{K}$ el conjunto de pendientes observadas en las $K$ ventanas de \num{30} min que cubren el ciclo de validación, y $\{\hat{y}_k\}_{k=1}^{K}$ las predicciones correspondientes del sistema híbrido. Se definen las siguientes métricas:

\paragraph{Coeficiente de determinación ($R^2$).}
Mide la fracción de varianza explicada por el modelo:
\begin{equation}
    \label{eq:cap4_r2_definicion}
    R^{2} = 1 - \frac{\sum_{k=1}^{K}\bigl(y_k - \hat{y}_k\bigr)^{2}}
    {\sum_{k=1}^{K}\bigl(y_k - \bar{y}\bigr)^{2}},
\end{equation}
donde $\bar{y}$ es la media de las observaciones. El criterio de éxito establecido en el marco metodológico exige $R^{2} \geq \num{0.70}$, indicando que al menos el \num{70}\% de la variabilidad en la tasa de emisión es capturada por la dinámica bioenergética del DEB y la corrección del GP. Un valor inferior sugeriría que el modelo mecanicista omite forzantes dominantes o que el estimador de biomasa no logra reconstruir la evolución larval \parencite{eriksenDynamicModellingFeed2022}.

\paragraph{Error porcentual absoluto medio (MAPE).}
Cuantifica la magnitud relativa del error de predicción:
\begin{equation}
    \label{eq:cap4_mape_definicion}
    \mathrm{MAPE} = \frac{100}{K}\sum_{k=1}^{K}
    \left|\frac{y_k - \hat{y}_k}{y_k}\right|.
\end{equation}
El objetivo global del sistema, conforme a la hipótesis H3, es alcanzar $\mathrm{MAPE} < \num{10}\%$ respecto a las mediciones de referencia (cromatografía de gases o sensor NDIR calibrado). Este umbral se selecciona porque representa la tolerancia aceptable para reportes de inventario de carbono en estándares MRV de segundo nivel (Tier 2) del IPCC, donde los factores de emisión derivados de modelos pueden sustituir mediciones directas si su incertidumbre se documenta rigurosamente \parencite{rossiEstimatingDynamicsGreenhouse2024}.

\paragraph{Raíz del error cuadrático medio (RMSE).}
Proporciona una medida de error en las mismas unidades que la variable objetivo (ppm/min):
\begin{equation}
    \label{eq:cap4_rmse_definicion}
    \mathrm{RMSE} = \sqrt{\frac{1}{K}\sum_{k=1}^{K}\bigl(y_k - \hat{y}_k\bigr)^{2}}.
\end{equation}

La comparación contra el modelo lineal de referencia se realiza mediante el cociente
\begin{equation}
    \label{eq:cap4_ratio_rmse}
    \rho_{\mathrm{RMSE}} = \frac{\mathrm{RMSE}_{\mathrm{híbrido}}}
    {\mathrm{RMSE}_{\mathrm{lineal}}},
\end{equation}
donde $\rho_{\mathrm{RMSE}} < 1$ indica superioridad del sistema DEB+GP. Se considera una mejora sustantiva si $\rho_{\mathrm{RMSE}} \leq \num{0.85}$, es decir, una reducción del error predictivo de al menos el \num{15}\% respecto a la extrapolación empírica simple. Este margen justifica el costo computacional y de implementación del estimador GP y la instrumentación asociada \parencite{biagiDevelopmentMachineLearningbased2024}.

Adicionalmente, se reporta el \emph{error máximo absoluto} (MaxAE) y el \emph{error porcentual absoluto máximo} (MaxAPE), identificando las ventanas temporales donde el sistema exhibe su mayor desviación. Estos picos de error suelen coincidir con eventos transitorios (post-alimentación, perturbaciones térmicas por apertura del reactor) y su localización temporal informa sobre las limitaciones estructurales del modelo.

\subsection{Análisis de residuos y detección de eventos de anaerobiosis}
\label{subsec:cap4_residuos_anaerobiosis}

Más allá de las métricas agregadas de regresión, el análisis de los residuos $\delta(t_k)$ definidos en la \cref{eq:cap4_residuo_operativo} proporciona diagnóstico sobre la validez del supuesto de aerobiosis y la calidad del ajuste. Un modelo bien especificado debe producir residuos que se comporten como ruido blanco gaussiano; desviaciones sistemáticas indican estructura no capturada.

\paragraph{Autocorrelación temporal.}
Se calcula la función de autocorrelación de los residuos hasta un desfase de \num{5} ventanas (\num{2.5} h):
\begin{equation}
    \label{eq:cap4_autocorrelacion}
    r_{\delta}(\tau) = \frac{\sum_{k=1}^{K-\tau}\bigl(\delta_k - \bar{\delta}\bigr)
    \bigl(\delta_{k+\tau} - \bar{\delta}\bigr)}
    {\sum_{k=1}^{K}\bigl(\delta_k - \bar{\delta}\bigr)^{2}}.
\end{equation}
Valores de $|r_{\delta}(\tau)| > \num{0.3}$ para $\tau \geq 1$ sugieren que el modelo subestima la persistencia temporal del proceso —por ejemplo, inercia térmica del sustrato o efectos de memoria en el metabolismo larval— e indican la necesidad de incluir términos de retardo o variables de estado adicionales en el DEB \parencite{eriksenDynamicModellingFeed2022}.

\paragraph{Normalidad y heterocedasticidad.}
Mediante el test de Shapiro-Wilk se evalúa la hipótesis de normalidad de los residuos. El rechazo de la normalidad ($p < \num{0.05}$) implica que la distribución predictiva del GP está mal calibrada o que existen outliers sistemáticos. La heterocedasticidad se inspecciona visualmente mediante un gráfico de residuos $\delta_k$ versus valores predichos $\hat{y}_k$: una banda de dispersión que se ensancha con $\hat{y}_k$ indica que la varianza del ruido instrumental $\sigma^{2}_{\mathrm{inst}}$ no es constante, violando uno de los supuestos de la función de pérdida \eqref{eq:cap4_funcion_perdida}.

\paragraph{Detección de eventos de anaerobiosis.}
La hipótesis H4 postula que el sistema actúa como alerta temprana de condiciones anaeróbicas mediante la lectura de \ce{CH4}. La validación de esta capacidad se formaliza como un problema de clasificación binaria donde:
\begin{itemize}
    \item \emph{Positivo verdadero} ($\mathrm{TP}$): el sensor detecta $C_{\mathrm{CH_4}} > \num{100}$ ppm y el operador confirma anaerobiosis (sustrato saturado, olor a sulfhídrico, coloración oscura).
    \item \emph{Falso positivo} ($\mathrm{FP}$): el sensor detecta \ce{CH4} pero no hay evidencia operativa de anaerobiosis (ruido transitorio del sensor, calibración deficiente).
    \item \emph{Negativo verdadero} ($\mathrm{TN}$): ausencia de \ce{CH4} y condiciones aeróbicas normales.
    \item \emph{Falso negativo} ($\mathrm{FN}$): presencia de anaerobiosis no detectada por el sensor (fallo instrumental o umbral mal calibrado).
\end{itemize}

Las métricas de clasificación se definen:
\begin{align}
    \label{eq:cap4_recall}
    \mathrm{Recall} &= \frac{\mathrm{TP}}{\mathrm{TP} + \mathrm{FN}}, \\
    \label{eq:cap4_precision}
    \mathrm{Precisión} &= \frac{\mathrm{TP}}{\mathrm{TP} + \mathrm{FP}}, \\
    \label{eq:cap4_f1}
    F_1 &= 2 \cdot \frac{\mathrm{Precisión} \cdot \mathrm{Recall}}
    {\mathrm{Precisión} + \mathrm{Recall}}.
\end{align}

El criterio de éxito para H4 exige $\mathrm{Recall} > \num{0.80}$: el sistema debe capturar al menos el \num{80}\% de los eventos de anaerobiosis reales. Se prioriza el recall sobre la precisión porque, en gestión ambiental, el costo de omitir una condición anaeróbica (emisión fugitiva de \ce{CH4}, degradación del sustrato, mortalidad larval) supera el costo de una falsa alarma que puede resolverse con una inspección visual \parencite{rossiEstimatingDynamicsGreenhouse2024}. La precisión, no obstante, se monitoriza para evitar la fatiga del operador; un valor aceptable se sitúa por encima del \num{60}\%.

\paragraph{Integración en el semáforo de riesgo.}
Los resultados del análisis de residuos y la clasificación de \ce{CH4} se sintetizan en el dashboard operativo mediante un indicador semafórico de tres niveles, coherente con el protocolo descrito en el \cref{chap:marco_metodologico}:

\begin{description}
    \item[Verde ($\mathcal{A}=0$, $|\delta| < \num{50}$ ppm/min)] El sistema opera dentro del régimen aeróbico predicho. El modelo DEB explica adecuadamente la dinámica de emisiones y el GP proporciona una estimación de biomasa coherente. Se acumulan créditos de carbono en el inventario dinámico.
    
    \item[Ámbar ($\mathcal{A}=0$, $|\delta| \geq \num{50}$ ppm/min)] Discrepancia significativa modelo-sensor sin presencia de \ce{CH4}. Sugiere una desviación metabólica no anaeróbica: posible inhibición por sustrato, estrés térmico, o error en la estimación del GP. Se recomienda revisión de la curva de crecimiento larval y verificación de la calibración del sensor NDIR.
    
    \item[Rojo ($\mathcal{A}=1$)] Detección confirmada de \ce{CH4}. Se activa la alerta de anaerobiosis y se instruye intervención inmediata: volteo del sustrato, ajuste de aireación o reducción de la densidad de alimentación. El residuo $\delta(t_k)$ en estas ventanas se excluye de la función de pérdida para calibración, conforme a la \cref{subsec:cap4_tratamiento_ch4}.
\end{description}

La validación predictiva, materializada en las métricas $R^2$, MAPE, RMSE y Recall, constituye la evidencia empírica que habilita el uso operativo del sistema. En la siguiente sección se presentan los resultados cuantitativos obtenidos durante los ciclos experimentales de validación.
